\chapter{Introduction}
The complex and inherently volatile nature of financial markets requires adaptive, data-driven strategies to optimize returns, coupled with effective risk management. Traditional technicals, like MACD or RSI, have provided the framework for trend and reversal identification for years. But static assumptions constrain their flexibility with respect to contemporary market dynamics. In this study, we provide a multi-dimensional research evolution— covering the historical use of technical indicators, deep learning (DL) predictions, and finally Deep Reinforcement Learning (DRL)— to align theoretical framework with actual trading performance. The phases provide distinct insights that lead to the FinDRL (Financial Deep Reinforcement Learning) framework to create a scalable tool for optimizing trading strategies. The exploration begins with an in-depth analysis of the 20 technical trading strategies of the FinDRL-* family implemented in PineScript V6 on TradingView Platform. Applying daily NVIDIA (NVDA) data from 2007 until 2025, strategies like FinDRL - Momentum (profit of 38.62 USD, 1,065 trades) and FinDRL - Volume Spike (profit factor 2.469) surpassed others like FinDRL - Bollinger Bands profit at -0.78 USD, exposing discrepancies in robustness. Backtested with a Stochastic RSI filter, these strategies — including trend-following, breakout, mean-reversion, and volume-based approaches — provide a systematic basis for quantitative evaluation, adaptive parameter tuning, and risk management. Based on 500+ strategies over 2 years distilled in the past year, this foundational work seeks to connect the big advances in mathematics to actionable trading, even while showing some pitfalls of static indicators to changes in the system they’re watching.

In the second phase, to fill in these gaps, we propose a DL-based framework for stock market prediction, "Stock Market Prediction: Forecasting and Decision-Making using GRU, CNN, and LSTM Networks", that predicts S\&P 500 (GSPC) prices. Using Gate Recurrent Units (GRU), Convolutional Neural Networks (CNN) and Long-Short Term Memory (LSTM) models, this method is able to extract non-linear patterns and long-term dependencies from raw price data. Using dynamic backtesting (with a starting balance of \$5,000 and a transaction cost of \$0.01), it addresses data non-stationarity with MinMaxScaler, optimizes lookback periods (between 10 and 120 days), and uses Keras Tuner for hyperparameter tuning. These ensemble predictions improve accuracy, and can be converted into trading signals. This framework increases adaptability with respect to the best technical indicators (e.g. Momentum), providing a practical approach for traders and tested in real-time using TradingView visualizations. It ends up with a DRL framework, FinDRL, combining deep neural networks (DNNs) for feature extraction and Deep Q-Networks (DQNs) for decision-making. FinnDRL utilizes a Gym-based trading environment, receives data enriched from a DataProcessor (e.g., Bollinger Bands, MACD), optimizes DNN hyperparameters with Bayesian optimization, and trains a neural network to learn the optimal buy/hold/sell actions in a way to yield a high cumulative reward. With a single-asset focus to start, it walks the line between immediate gains and longer-term strategy, outmaneuvering the rigid smart order routing of the traditional algorithms. Its modular architecture promotes replicability and extensibility, setting a basis for multi-asset expansions and further complex risk-adjusted measures.

In sum, this work contributes to quantitative finance by moving from indicators to DL-driven forecasting then DRL based decisions. It serves as a solid, flexible toolbox for practitioners and academics, catering to topics including real-time adaptability, risk assessment, and the merging of theory with practice—laying the groundwork for future dynamic innovations in financial market environments.

\section{Thesis Structure}
Key Concepts in Financial Deep Reinforcement Learning - This research unfolds across four chapters, each rooted in the concepts below.
\begin{enumerate}
    \item \textbf{Trading Strategies Architectural Design: A Comparative Analysis of 20 Technical Indicators, Using a Dynamic Backtesting, and Performance Evaluation on NVIDIA Stock} \\
    Concept: This basic study spark to test 20 technical indicators —Momentum, Bollinger Bands, Volume Spike and others—that I coded in PineScript and applied to NVDA daily data (2007–2025) With a Stochastic RSI filter, these strategies harness trends \& shifts in volume \& reversal. Dynamic backtesting evaluates performance as you would expect it in real life with a starting balance of \$10,000 and \$0.01 fees: Momentum benefits from 1,065 trades earning +\$38.62 and Parabolic SAR loses \$9.89 The architecture contrasts trend-following, mean-reversion and breakdown strategies, exposing strengths (e.g. Volume Spike’s 2.469 profit factor) and weaknesses (e.g. RSI’s idleness). Think of it as a scientific road map for an empirical study of NVDA’s wild swings.

    \item \textbf{Maximizing Market Returns Optimizing Technical Indicators: A Bayesian Optimization Approach} \\
    Concept: This method builds on the earlier indicator analysis using Optuna along with the Bayesian optimization process to optimize parameters of the 20 strategies for NVDA. Even 14-day lookback for momentum and 20-day window for Bollinger Bands get data-driven new leases on life, with profit vs drawdown balanced across 50 trials. A Gaussian Process dictates it, adapting to the cadence of NVDA-schwing—\$9.89 loss on 10-day Momentum for 2023 AI spikes or tamed Parabolic SAR. This approach enhances returns and flexibility, transforming static tools into proactive assets. It’s about accuracy, supporting the initial study’s findings with a scalable, automated optimization perspective.

    \item \textbf{Deep Learning for Stock Market Prediction: Forecasting and Decision-Making using GRU, CNN, and LSTM Networks} \\
    Concept: This article changes tides modeling S\&P500 (GSPC) prices using deep neural networks—GRU, CNN, and LSTM. GRU snags short-term jolts (e.g. 2024 chip releases), CNN finds local patterns (e.g. earnings dips), and LSTM tracks long-term trends (2020 GPU rises). Models trained on those indicators — optimized like Stochastic RSI — ensemble through a \texttt{CombinedModel}, hitting \$50 profit in backtests (\$5,000 start, \$0.01 cost) — the highest of any model beating Momentum’s \$38.62. Panama uses non-linear chaos on S\&P500 (GSPC) providing actionable signals with validations on TradingView. The idea is predictive power, using previous optimizations to connect raw data to real world decisions.
    
    \item \textbf{FinDRL: Financial Deep Reinforcement Learning - A Perfect Way to Manage Stock Trading Portfolio} \\
    Concept: FinDRL applies them all to trade the S\&P500 (GSPC) using deep reinforcement learning. A DNN in a \texttt{TradingEnv} extracts features [OHLCV, optimized indicators] and learns how to take actions by Q-learning for buy, hold, sell across 100 episodes using a DQN. It gets 15–20\% growth and a 1.2 Sharpe ratio, buying troughs of 2022 and selling tops of 2023. A synthesis of its earlier efforts—indicators/crunchers, Bayesian tweaks, DNN forecasts—it adapts where static strategies stumble, melding short-term and long-term compound objectives. There it is: an ever-adapting,self-optimizing engine capable of taming S\&P500 (GSPC)'s volatility into a deployable, scalable portfolio management tool.
    
\end{enumerate}